% !TEX root = ../report.tex

\chapter{Технологическая часть}

В данном разделе описывается реализация базы данных и программного интерфейса доступа к ней. Рассматриваются средства реализации, порядок применения миграций, реализованные ограничения целостности, роли СУБД, триггеры, хранимая функция, Redis-кэширование и REST-интерфейс.

\section{Выбор средств реализации}

Для реализации базы данных используется PostgreSQL 16. PostgreSQL относится к реляционным СУБД и поддерживает транзакции, внешние ключи, ограничения целостности, частичные индексы, пользовательские роли, триггеры и хранимые функции на языке PL/pgSQL~\cite{postgresql_doc}. Эти возможности соответствуют требованиям проектируемой системы: данные маркетплейса имеют выраженную табличную структуру, а часть правил целостности должна выполняться непосредственно на уровне базы данных.

Выбор СУБД выполнен по критериям, связанным с реализуемой схемой. В качестве альтернатив рассматривались PostgreSQL, MySQL и SQLite~\cite{postgresql_doc,mysql_doc,sqlite_doc}. Сравнение приведено в таблице~\ref{tbl:dbms_compare}.

\begin{table}[H]
    \caption{Сравнение СУБД по требованиям реализации}
    \label{tbl:dbms_compare}
    \begin{tabular}{|p{0.36\textwidth}|c|c|c|}
        \hline
        \textbf{Критерий} & \textbf{PostgreSQL} & \textbf{MySQL} & \textbf{SQLite} \\ \hline
        Серверный многопользовательский режим & да & да & нет \\ \hline
        Роли и привилегии на уровне СУБД & да & да & нет \\ \hline
        Хранимые функции с табличным результатом & да & да & нет \\ \hline
        Частичные индексы & да & нет & да \\ \hline
        Статическая типизация столбцов (\texttt{NUMERIC}, \texttt{TIMESTAMP}) & да & да & нет \\ \hline
        Поддержка внешних ключей и транзакций & да & да & да \\ \hline
    \end{tabular}
\end{table}

По результатам сравнения PostgreSQL покрывает все перечисленные требования. В реализации используются роли СУБД, частичные индексы, триггеры и функция, возвращающая табличный результат. Поэтому выбрана СУБД PostgreSQL.

Серверное приложение реализовано на языке Go. Язык используется для реализации HTTP-интерфейса, бизнес-операций, транзакционных сценариев и взаимодействия с PostgreSQL и Redis. В проекте используется Go версии 1.25.0 (директива \texttt{go 1.25.0} в \texttt{go.mod})~\cite{go_doc}. Для работы с PostgreSQL используется драйвер \texttt{pgx/v5}, для маршрутизации HTTP-запросов --- библиотека \texttt{chi}, для построения SQL-запросов --- библиотека \texttt{squirrel}.

Миграции схемы выполняются утилитой \texttt{goose}~\cite{goose_doc}. Каждая миграция содержит секции \texttt{Up} и \texttt{Down}, что позволяет воспроизводимо применять и откатывать изменения схемы. Такой подход фиксирует порядок создания таблиц, ограничений, индексов, ролей, триггеров и функций.

Redis используется как in-memory хранилище для кэша часто читаемых данных и для хранения списка отозванных JWT-токенов~\cite{redis_doc}. PostgreSQL остаётся источником истины; Redis используется только как дополнительный слой ускорения чтения и проверки сессий.

В реализации используются PostgreSQL~16.11, Redis~8.6 и Go~1.25.0. Для работы с PostgreSQL применяется драйвер \texttt{pgx/v5}, для маршрутизации HTTP-запросов --- \texttt{chi}, для построения параметризованных SQL-запросов --- \texttt{squirrel}, для применения миграций --- \texttt{goose}.

\section{Реализация базы данных}

Схема базы данных создаётся набором миграций, расположенных в каталоге \texttt{migrations}. Состав миграций приведён в таблице~\ref{tbl:tech_migrations}. Полные SQL-тексты миграций приведены в приложении~А.

\begin{table}[H]
    \caption{Миграции базы данных}
    \label{tbl:tech_migrations}
    \begin{tabular}{|p{0.22\textwidth}|p{0.70\textwidth}|}
        \hline
        \textbf{Группа миграций} & \textbf{Назначение} \\ \hline
        \texttt{001}, \texttt{006}, \texttt{009}--\texttt{012} & создание таблиц и добавление ограничений предметной области: связь заказа с продавцом, резерв товара, адрес доставки, инварианты корзины и изменение уникальности телефона \\ \hline
        \texttt{002} & создание индексов для выборок товаров, заказов, отзывов, категорий и активных записей \\ \hline
        \texttt{003}, \texttt{008} & создание ролей PostgreSQL, назначение привилегий и добавление системных учётных записей \\ \hline
        \texttt{004}, \texttt{005}, \texttt{007} & создание триггера журнала изменения цен, функции статистики продавца и триггера пересчёта рейтингов \\ \hline
    \end{tabular}
\end{table}

Итоговая схема содержит десять таблиц. Таблицы и ограничения были описаны в конструкторской части в таблицах~\ref{tbl:t_users}--\ref{tbl:t_price_history}. В технологической реализации сохраняются следующие группы ограничений:

\begin{itemize}
    \item ограничения домена значений: положительная цена товара, неотрицательное количество товара, допустимый диапазон оценки отзыва, фиксированный набор ролей и статусов заказа;
    \item ограничения ссылочной целостности: связи товара с продавцом, заказа с покупателем, позиции заказа с заказом и товаром, отзыва с пользователем и товаром;
    \item ограничения уникальности: электронная почта пользователя, профиль продавца на пользователя, имя категории, пара заказа и товара в позиции заказа, пара пользователя и товара в отзыве;
    \item частичный уникальный индекс для ограничения одной draft-корзины на покупателя;
    \item проверка \texttt{reserved\_quantity <= stock\_quantity}, исключающая резервирование количества больше доступного остатка.
\end{itemize}

Для мягкого удаления пользователей и товаров используются столбцы \texttt{deleted\_at}. Поиск активных пользователей и товаров поддерживается частичными индексами \texttt{idx\_users\_id\_not\_deleted} и \texttt{idx\_products\_id\_not\_deleted}. Эти индексы включают только строки, у которых \texttt{deleted\_at IS NULL}.

\section{Реализация ролевой модели}

На уровне PostgreSQL реализованы четыре роли:

\begin{itemize}
    \item \texttt{marketplace\_buyer};
    \item \texttt{marketplace\_seller};
    \item \texttt{marketplace\_admin};
    \item \texttt{marketplace\_analyst}.
\end{itemize}

Роли создаются миграцией \texttt{003\_create\_roles.sql}. Привилегии назначаются командами \texttt{GRANT}.

Роль покупателя получает права чтения каталога и категорий, а также права изменения адресов, отзывов, заказов и позиций заказа. Роль продавца получает права управления товарами и чтения заказов, связанных с его товарами. Роль администратора получает полный доступ к таблицам и последовательностям. Роль аналитика получает доступ только на чтение.

Роли PostgreSQL задают технический уровень доступа к объектам БД. Правила принадлежности отдельных строк проверяются в приложении при выполнении пользовательских операций. Например, покупатель работает только со своими адресами, заказами и отзывами, а продавец --- со своим профилем, своими товарами и заказами, относящимися к этим товарам.

\section{Реализация триггеров и хранимой функции}

В базе данных реализованы два триггера и одна хранимая функция. Их назначение соответствует проектным решениям, описанным в конструкторской части.

Триггер \texttt{catch\_price\_change} создаётся миграцией \texttt{004}. Он срабатывает после изменения столбца \texttt{price} в таблице \texttt{products}. Условие \texttt{OLD.price IS DISTINCT FROM NEW.price} исключает запись в журнал при отсутствии фактического изменения цены. При срабатывании триггера функция аудита добавляет строку в таблицу \texttt{product\_price\_history}.

Триггер \texttt{update\_ratings\_on\_review} создаётся миграцией \texttt{007}. Он срабатывает после добавления отзыва, изменения оценки или удаления отзыва. Функция пересчёта рейтингов обновляет рейтинг товара как среднее значение оценок по этому товару. Затем пересчитывается рейтинг продавца как среднее значение рейтингов его активных товаров.

Хранимая функция \texttt{get\_seller\_statistics} создаётся миграцией \texttt{005}. Функция принимает идентификатор продавца и границы периода. В результат входят количество заказов, суммарная выручка, средний чек и наименование самого продаваемого товара. В расчёт включаются только заказы со статусами \texttt{paid}, \texttt{shipped} и \texttt{delivered}. Цена берётся из столбца \texttt{price\_at\_purchase}, поэтому изменение текущей цены товара не меняет статистику уже оформленных заказов.

\section{Тестирование хранимой функции}

Для проверки функции \texttt{get\_seller\_statistics} выделены классы эквивалентности входных данных. Классы приведены в таблице~\ref{tbl:stats_equiv}.

\begin{table}[H]
    \caption{Классы эквивалентности для тестирования функции \texttt{get\_seller\_statistics}}
    \label{tbl:stats_equiv}
    \begin{tabular}{|p{0.08\textwidth}|p{0.35\textwidth}|p{0.45\textwidth}|}
        \hline
        \textbf{Класс} & \textbf{Условие} & \textbf{Ожидаемый результат} \\ \hline
        K1 & у продавца нет завершённых продаж в заданном периоде & \texttt{total\_orders = 0}; агрегаты выручки, среднего чека и лидирующего товара не формируются \\ \hline
        K2 & у продавца есть один оплаченный заказ с одной позицией & количество заказов равно 1; выручка равна произведению количества на \texttt{price\_at\_purchase}; средний чек равен выручке \\ \hline
        K3 & один заказ содержит несколько позиций товаров продавца & заказ учитывается один раз; выручка суммируется по всем позициям \\ \hline
        K4 & в период попадают заказы со статусами \texttt{draft}, \texttt{pending} или \texttt{cancelled} & такие заказы исключаются из расчёта \\ \hline
        K5 & часть заказов находится вне заданного временного интервала & заказы вне интервала не влияют на агрегаты \\ \hline
        K6 & продано несколько товаров с разным количеством единиц & \texttt{top\_product\_name} соответствует товару с максимальной суммой \texttt{quantity} \\ \hline
        K7 & два или более товаров продавца имеют одинаковую максимальную суммарную проданную величину & возвращается один из лидирующих товаров; при равенстве сумм выбор конкретного товара не доопределяется и не влияет на остальные агрегаты \\ \hline
    \end{tabular}
\end{table}

Проверка должна выполняться на тестовых данных, содержащих заказы разных статусов, несколько товаров одного продавца и заказы за пределами выбранного периода. Такой набор данных позволяет проверить фильтрацию по статусу, фильтрацию по дате, подсчёт уникальных заказов и выбор лидирующего товара.

\section{Кэширование с использованием Redis}

Для часто повторяемых чтений используется стратегия Cache-Aside. При попадании в кэш приложение возвращает данные из Redis; при промахе читает данные из PostgreSQL, записывает их в Redis с TTL и возвращает результат. PostgreSQL остаётся источником истины. Redis не используется для хранения корзин, заказов, платежей, профилей и отчётов.

Схема взаимодействия PostgreSQL и Redis приведена на рисунке~\ref{fig:redis_interaction}.

\begin{figure}[H]
    \centering
    \begin{tikzpicture}[>=stealth, font=\small]
        \node[draw, rectangle, align=center, minimum width=3.4cm, minimum height=1cm] (client) at (0, 3.2) {HTTP-клиент};
        \node[draw, rectangle, align=center, minimum width=4.2cm, minimum height=1cm] (app) at (0, 1) {Go-приложение\\service/repository};
        \node[draw, rectangle, align=center, minimum width=3.2cm, minimum height=1cm] (redis) at (8.1, 1) {Redis\\кэш};
        \node[draw, rectangle, align=center, minimum width=3.6cm, minimum height=1cm] (pg) at (0, -2.1) {PostgreSQL\\источник истины};

        \draw[->] (client) -- node[right] {1. HTTP} (app);
        \draw[->] ([yshift=0.2cm]app.east) -- node[above] {2. GET} ([yshift=0.2cm]redis.west);
        \draw[->] ([yshift=-0.2cm]redis.west) -- node[below] {3. hit/miss} ([yshift=-0.2cm]app.east);
        \draw[->] (app.north east) to[out=35,in=145] node[above] {6. SET TTL} (redis.north west);
        \draw[->, dashed] (app.south east) to[out=-35,in=-145] node[below] {инвалидация: DEL/SCAN} (redis.south west);

        \draw[->] ([xshift=-0.45cm]app.south) -- node[left] {4. SQL} ([xshift=-0.45cm]pg.north);
        \draw[->] ([xshift=0.45cm]pg.north) -- node[right] {5. данные} ([xshift=0.45cm]app.south);
    \end{tikzpicture}
    \caption{Схема взаимодействия PostgreSQL и Redis}
    \label{fig:redis_interaction}
\end{figure}

Сплошные стрелки (1--6) образуют путь чтения по схеме Cache-Aside. Пунктирная стрелка относится к инвалидации ключей после изменения данных.

Ключи Redis приведены в таблице~\ref{tbl:redis_keys}.

\begin{table}[H]
    \caption{Ключи Redis}
    \label{tbl:redis_keys}
    {\small
    \begin{tabular}{|p{0.38\textwidth}|p{0.16\textwidth}|p{0.34\textwidth}|}
        \hline
        \textbf{Ключ} & \textbf{TTL} & \textbf{Назначение} \\ \hline
        \texttt{products:\{id\}} & \texttt{product\_ttl} & карточка товара \\ \hline
        \texttt{products:catalog:\{query\}} & \texttt{catalog\_ttl} & страница каталога с параметрами фильтрации, сортировки и пагинации \\ \hline
        \texttt{categories:list:\{query\}} & \texttt{category\_ttl} & список категорий \\ \hline
        \makecell[l]{\texttt{products:\{id\}:reviews:}\\\texttt{limit=\{l\}\&page=\{p\}}} & \texttt{review\_ttl} & страница отзывов товара \\ \hline
        \texttt{sessions:\{jti\}} & срок действия JWT & список отозванных токенов \\ \hline
    \end{tabular}
    }
\end{table}

Инвалидация выполняется после успешной фиксации изменений в PostgreSQL. Изменение товара, категории или отзыва очищает связанные ключи кэша. Redis также используется для хранения списка отозванных JWT-токенов.

\section{Интерфейс доступа к базе данных}

Доступ к данным реализован через JSON REST API с префиксом \texttt{/api/v1}. Интерфейс включает группы маршрутов авторизации, профиля пользователя, каталога, товаров продавца, корзины и заказов, оплаты, отзывов и администрирования. HTTP-семантика методов соответствует стандартному назначению операций чтения, создания, обновления и удаления ресурсов~\cite{http_semantics}. Для защищённых маршрутов используется JWT, передаваемый в заголовке \texttt{Authorization: Bearer <token>}; ошибки возвращаются в виде объекта \texttt{\{"error": "message"\}}.

Помимо JSON API, реализован серверный Web-интерфейс на основе \texttt{html/template}. Оба интерфейса используют общий слой операций и работают с одной схемой базы данных.

\section*{Вывод}
\addcontentsline{toc}{section}{Вывод}

Схема базы данных реализована на PostgreSQL набором миграций. Ограничения целостности обеспечиваются на уровне СУБД, разграничение доступа --- ролями PostgreSQL и проверками принадлежности записей в приложении. Аудит изменения цен и пересчёт рейтингов вынесены в триггеры, агрегатная статистика продавца --- в хранимую функцию. Redis используется для кэширования повторных чтений по схеме Cache-Aside и хранения отозванных токенов.
